\section{Background}
\label{sec:background}

\subsection{VFS Name Resolution}

The Linux VFS resolves a pathname by walking directory components, consulting
directory entries and inodes, and then dispatching ordinary file operations to
the selected filesystem~\cite{linux_vfs_docs,linux_path_lookup}. Lookup selects the object on which later file
semantics operate. File semantics, including permission checks, file reads and
writes, page-cache behavior, persistence, and consistency, remain
responsibilities of the VFS and lower filesystem. \namei builds on this
separation. It changes path-view selection during lookup and readdir, but it
does not replace the file operations that run after an object is selected.

\subsection{Programmable Kernel Policy}

eBPF lets the kernel run verified, bounded programs at explicit attachment
points~\cite{linux_bpf_docs,linux_bpf_verifier}. The \code{sched_ext} scheduler class is a useful precedent because it
lets BPF choose scheduling policy while the kernel retains scheduler
machinery~\cite{linux_sched_ext}. \namei uses the same policy-versus-ownership
separation for VFS name resolution. BPF supplies a bounded path-view decision,
and the kernel retains the machinery that owns filesystem objects. The analogy
concerns the ownership split rather than the scheduler interface.

\subsection{Namespace Construction And Filesystem Services}

Linux already offers mechanisms for alternate filesystem views. Bind mounts,
OverlayFS, projected volumes, and copied trees construct or materialize a
namespace before the application uses it~\cite{linux_mount,linux_overlayfs,kubernetes_projected_volumes}.
FUSE and stackable or custom filesystems provide a broader programmable
filesystem boundary, often by placing a service or filesystem implementation on
the operation path~\cite{linux_fuse,wrapfs,bento}. These mechanisms define the neighboring
boundaries for the path-view policy studied in the rest of the paper.
